% !TEX root = ../main.tex
%-----------------------------------------------------------------------
%  PART I -- the vertical-diffusivity perturbation ensemble.
%  \input by ../main.tex -- not compiled on its own.
%-----------------------------------------------------------------------
\part{The vertical-diffusivity perturbation ensemble}

\begin{callout}
\textbf{Status.} This part is written as the plan it was; the ensemble it
describes was executed in full on 2026-08-28/29 and analysed on 2026-08-30.
Section~\ref{sec:results} holds the results, the answer to the question below
(a strong response --- $\kv$ is a required surrogate input), and the revised
next steps. The design prose here is kept as the record of what was decided
and why.
\end{callout}

%-----------------------------------------------------------------------
\section{Objective and context}
\label{sec:objective}
%-----------------------------------------------------------------------

\subsection{The question}

Adjoint sensitivity maps are computed at one set of parameter values. Whether
those maps remain valid when the parameters change is not known for this
configuration, and it is not a question a single adjoint run can answer.

The experiment asks directly: \textbf{do the adjoint sensitivity patterns depend
on the model's vertical mixing?} The answer has consequences in two directions.

\begin{itemize}[nosep]
  \item If the patterns are largely insensitive to $\kv$ across a wide range,
        then published sensitivity maps are robust to the mixing uncertainty that
        any ocean model carries. That is a useful result in its own right,
        independent of any machine learning.
  \item If the patterns change substantially with $\kv$, then any surrogate
        trained to predict them must take mixing as an input, and its training
        set must span the mixing range.
\end{itemize}

\subsection{A second result: validating the adjoint}
\label{sec:validating}

The ensemble also validates the adjoint itself. The adjoint returns
$\partial J/\partial\kv$ at the reference point; the ensemble returns
$\Delta J/\Delta\kv$ by finite difference across members. Where the two diverge
marks the perturbation size beyond which the linearisation stops holding.

This matters more than it may appear, because \textbf{the adjoint is not
currently validated by anything else.} MITgcm's built-in gradient check runs on
every adjoint job in this configuration, but it perturbs a point near the
southern boundary while the cost section sits at 26\dg N. Sensitivity there is of
order $10^{-9}$ against a field maximum of order $10^{-2}$, so the check measures
run-to-run noise rather than the perturbation, and it fails by roughly eight
orders of magnitude --- as it always has. Appendix~\ref{app:grdchk} gives the
detail. The ensemble's finite-difference comparison, taken at points where the
sensitivity is large, would be the first meaningful check this configuration has
had.

\subsection{What Part~II depends on}
\label{sec:bridge}

Part~II of this document sets out a neural-network surrogate that predicts
adjoint sensitivity fields directly instead of computing them. The two parts are
sequential rather than parallel:

\begin{itemize}[nosep]
  \item \textbf{This ensemble is a scoping study.} Its primary output is a design
        answer --- whether $\kv$ must be an input to the surrogate at all. The
        training pairs it produces are a by-product.
  \item \textbf{Its result determines the surrogate's design.} A flat response
        redirects the surrogate towards ocean state and cost functional as its
        varying inputs. A strong response makes mixing a required input and
        justifies the larger parameter ensemble of
        Section~\ref{sec:trainingdata}.
\end{itemize}

Eight $\kv$ values, the seven members plus the reference, are ample for the first
question and thin for the second. Part~II therefore calls for a $200$--$400$
member ensemble later; Part~I does not propose one now.

Three of the open questions in Section~\ref{sec:questions} --- the spatial
structure of the perturbation, the surrogate framing, and the scope of the cost
function --- bear on both parts and should be settled once for both.

%-----------------------------------------------------------------------
\section{Experiment design}
%-----------------------------------------------------------------------

Each member follows the same sequence:

\begin{enumerate}[nosep]
  \item Start from the spin-up state at \textbf{year 2170}.
  \item Perturb \textbf{vertical diffusivity only}. GM and KPP are already
        disabled.
  \item Run \textbf{10 years forward} to re-equilibrate, arriving at
        \textbf{year 2180}.
  \item Run a \textbf{5-year adjoint over 2180 $\rightarrow$ 2185}, identical in
        configuration to the reference run.
  \item Compare the resulting sensitivity patterns against the reference.
\end{enumerate}

\subsection{Members}

Seven perturbations, geometric in powers of two, spanning a factor of 128 around
the reference value of $1.2\times10^{-5}\,\mathrm{m^{2}\,s^{-1}}$:

\begin{center}
\footnotesize
\begin{tabular}{l*{8}{r}}
\toprule
$\kv$ ($\mathrm{m^{2}\,s^{-1}}$)
 & 3.0e$-$6 & 6.0e$-$6 & \textbf{1.2e$-$5} & 2.4e$-$5
 & 4.8e$-$5 & 9.6e$-$5 & 1.9e$-$4 & 3.8e$-$4\\
$\times$ reference
 & 0.25 & 0.5 & \textbf{1 (done)} & 2 & 4 & 8 & 16 & 32\\
\bottomrule
\end{tabular}
\end{center}

All values lie inside the admissible range already declared for this control,
$1\times10^{-6}$ to $5\times10^{-4}\,\mathrm{m^{2}\,s^{-1}}$. The spacing is
logarithmic because the effect of mixing on stratification scales closer to
$\log\kv$ than to $\kv$, and the reference sits low in the admissible window,
leaving most of the range above it.

\subsection{Why the perturbation leg starts at 2170}
\label{sec:leg}

Placing the perturbation leg \emph{before} the reference window makes the
completed run a valid control at no additional cost:

\begin{center}
\footnotesize
\begin{verbatim}
member    :  state @ 2170 --10 yr, perturbed k--> state @ 2180 --> adjoint 2180-2185
reference :  state @ 2170 --10 yr, reference k--> state @ 2180 --> adjoint 2180-2185
\end{verbatim}
\end{center}

The spin-up ran through 2170 $\rightarrow$ 2180 continuously, so its own
trajectory is the unperturbed counterpart of each member's leg. Both branches
begin from an identical state, run for the same duration, and are evaluated over
the same window, which leaves diffusivity as the only difference between them.

Re-equilibrating forward from 2180 instead would place member windows at
2190 $\rightarrow$ 2195, conflating the $\kv$ change with a decade of internal
variability. The AMOC wanders by several tenths of a Sverdrup over such an
interval, which is comparable to the signal being sought.

The design therefore requires \textbf{seven new runs, not eight}. One caveat is
inherent to any finite re-equilibration: ten years may not fully equilibrate a
member. That is open question 4 in Section~\ref{sec:questions}, and the pilot
member tests it.

%-----------------------------------------------------------------------
\section{Compute and storage}
\label{sec:compute}
%-----------------------------------------------------------------------

Both halves of a member have been measured on completed runs at this exact
configuration and rank count, so the figures below are not extrapolations.

\begin{center}
\begin{tabularx}{\textwidth}{@{}Y{1.35}Y{1.15}Y{1.0}Y{0.5}@{}}
\toprule
\textbf{Measured from} & \textbf{Work done} & \textbf{Wall time, 27 ranks} &
\textbf{Per step}\\
\midrule
200-year spin-up (job 28463) & $3{,}513{,}600$ forward steps & 32 h 21 m
  & \textbf{0.0332 s}\\
5-year reference adjoint (job 28486) & $87{,}840$ adjoint steps & 23 h 03 m
  & \textbf{0.945 s}\\
30-day adjoint (job 30948) & $1{,}440$ adjoint steps & 23 m 19 s
  & 0.972 s\\
\bottomrule
\end{tabularx}
\end{center}

The two adjoint measurements agree to 3\,\%, which is the useful cross-check:
cost is linear in window length, so a 5-year adjoint is predictable from a
20-minute one.

A member therefore costs \textbf{1.6 h forward plus 23.0 h adjoint, about
24.7 h}, and seven members about 173 h. That is the cost as currently configured,
and most of it is avoidable.

\subsection{The adjoint is doing four times the work it needs to}
\label{sec:grdchk-cost}

The ratio of adjoint to forward cost per timestep is \textbf{28.5}, far above the
3--5 that checkpointing alone would explain. The cause is that
\code{useGrdchk = .TRUE.} is set, so every adjoint job also runs the built-in
finite-difference gradient check and pays for the perturbed forward integrations
it requires. Instrumenting the 30-day run:

\begin{center}
\footnotesize
\begin{verbatim}
forward-step calls executed   18,622
timesteps actually integrated  1,440      -> 12.9x
   of which: gradient check   ~11,520      (4 check points, 2 perturbed runs each)
             adjoint proper    ~7,102      (forward sweep + checkpoint recomputation)
\end{verbatim}
\end{center}

The adjoint proper therefore costs about 4.9 forward-step equivalents per
timestep, which is within the 3--5 range checkpointing would predict. The excess
is entirely the gradient check.

Disabling that check for ensemble members should cut adjoint work by about
\textbf{2.6$\times$}, bringing a member to roughly
\textbf{1.6 h + 9 h $\approx$ 10.5 h} and the ensemble to \textbf{about 75 h}.
Nothing is lost by doing so: the check cannot pass at its current setting
(Section~\ref{sec:validating}), and the ensemble's own finite-difference
comparison is a stronger validation of the same property.

\begin{center}
\begin{tabularx}{\textwidth}{@{}Y{1.6}Y{0.6}Y{0.8}@{}}
\toprule
\textbf{Configuration} & \textbf{Per member} & \textbf{Seven members}\\
\midrule
As the reference run was configured & 24.7 h & \textbf{173 node-h}\\
With the gradient check disabled & $\sim$10.5 h & \textbf{$\sim$75 node-h}\\
Gradient check disabled, 4 runs packed per node & --- & \textbf{$\sim$19 node-h}\\
\bottomrule
\end{tabularx}
\end{center}

\subsection{Why 200 node-hours is requested}

The measurements come from the machine the reference run used; the target
architecture has not been benchmarked. Strongly perturbed members may require a
reduced timestep or additional solver iterations, and some of seven
never-before-integrated settings will need repeating. The 2.6$\times$ saving is
inferred from call counts rather than measured on a completed run, so it is not
yet banked. A short benchmark, well under one node-hour and the first step of the
sequence in Section~\ref{sec:sequence}, settles all of this.

\paragraph{Node packing.}
At 27 ranks per run, four runs fit on a 128-core node, an occupancy of 84\,\%.
One run per node would leave 79\,\% of each billed node idle, so packing can cut
charged node-hours by up to 4$\times$ subject to memory-bandwidth contention.
This is worth measuring in the same benchmark.

\paragraph{Storage.}
About 3.3 GB per run --- 1.7 GB of sensitivity output and 1.6 GB of forward
diagnostics --- giving roughly \textbf{23 GB} for the ensemble. Storage is not a
constraint. Output frequency is the lever: 5-day sampling is assumed, and finer
output scales linearly. One caution from the reference configuration: it retains
a restart every 30.5 days, and a 10-year member leg would write 120 such files,
about 2.8 GB, of which only the last is needed. Raising \code{pChkptFreq} for the
members saves roughly 20 GB across the ensemble at no cost.

%-----------------------------------------------------------------------
\section{Analysis plan}
%-----------------------------------------------------------------------

\subsection{Does the sensitivity depend on \texorpdfstring{$\kv$}{kappa}?}

The primary question. Each member is compared against the reference run, and
because both share the 2180 $\rightarrow$ 2185 window the comparison is direct:

\begin{itemize}[nosep]
  \item pattern correlation of each \code{ADJ*} field, member against reference;
  \item amplitude ratio;
  \item whether the sign structure survives.
\end{itemize}

These are plotted against $\log\kv$. A flat line means the sensitivities are
robust to mixing. A strong trend means $\kv$ must be an input to any surrogate.

If a noise floor is later needed to judge weak trends against, the cheapest
source is a second unperturbed run started from a different checkpoint; that
difference is pure internal variability at fixed $\kv$. It is not needed to
interpret a strong signal.

\subsection{Validating the adjoint against finite differences}
\label{sec:fd}

The adjoint gives $\partial J/\partial\kv$ at the reference point, and the
ensemble gives $\Delta J/\Delta\kv$ between members. For small perturbations the
two should agree; as $\Delta\kv$ grows they will diverge. The point of divergence
gives the range over which adjoint sensitivities can be trusted, and indicates
whether a surrogate must be nonlinear in $\kv$.

This extends the built-in gradient check from an infinitesimal perturbation to a
physically meaningful one, and does so at a point where the sensitivity is large
enough to measure.

\subsection{A caution before comparing across controls}

Every control currently carries unit weight
(\code{xx\_*\_weight = 'ones\_64b.bin'}), so the raw gradients are expressed in
each control's own units and \textbf{are not comparable across controls}. In the
reference run \code{adxx\_empmr} spans about $10^{4}$ and \code{adxx\_qsw} about
$10^{-6}$, which is an artefact of units rather than physics.

Any cross-control statement, and any attempt to feed several controls to a single
network, requires normalising each by a prior uncertainty or at least by a
characteristic scale.

%-----------------------------------------------------------------------
\section{Ensemble axes beyond vertical diffusivity}
%-----------------------------------------------------------------------

Diffusivity is the first axis, not the only one.

\begin{center}
\begin{tabularx}{\textwidth}{@{}Y{0.3}Y{1.8}Y{0.9}@{}}
\toprule
\textbf{Axis} & \textbf{Content} & \textbf{Relative cost}\\
\midrule
1 & Vertical diffusivity --- Part~I & as above\\
2 & \textbf{Ocean state sampled along the existing spin-up}
  & adjoint only, no forward leg\\
3 & Surface forcing amplitude, wind stress first & about one member each\\
--- & Initial-condition perturbations; additional cost functions & deferred\\
\bottomrule
\end{tabularx}
\end{center}

\paragraph{Axis 2 is nearly free and deserves emphasis.}
The spin-up wrote a restart state every 30.5 days and \textbf{all $2{,}400$
remain on disk}; this has been verified rather than assumed. Each is a
dynamically balanced ocean differing from its neighbours through internal
variability alone. An adjoint from each needs no forward integration, so it costs
about 90\,\% of a diffusivity member rather than the two-thirds assumed earlier
--- the forward leg turned out to be far cheaper than the adjoint. The saving is
smaller than first thought, but the argument stands: Axis 2 samples states the
model actually visits rather than states imposed by hand, $2{,}400$ candidates
are available today, and no new forward integration is required. It could run
concurrently with Axis 1.

\paragraph{Initial-condition perturbations are not worth a separate ensemble.}
They probe the same question as state sampling --- how sensitivity depends on
ocean state under fixed dynamics --- but from an unbalanced state. Surface
anomalies are damped within days by the roughly 6.5-day restoring while interior
anomalies persist, so what actually perturbs the run is whatever survives, which
is difficult to control.

Axis 3 would roughly double the request and would be raised separately after
Axis 1.

%-----------------------------------------------------------------------
\section{Open questions}
\label{sec:questions}
%-----------------------------------------------------------------------

These are the points on which reviewer input would be most useful.

\begin{enumerate}
  \item \textbf{Member count and spacing.} Seven members, logarithmically
        spaced. Is that the right trade against fewer members with longer
        re-equilibration?

  \item \textbf{Should the perturbation carry spatial structure?} The reference
        diffusivity is spatially uniform, so as proposed every member is also a
        constant field. A real ocean has elevated mixing over rough topography.
        This also bears directly on the surrogate: its stated value lies in
        handling \emph{field-valued} parameters with hundreds of thousands of
        components, but if every training field is a constant then $\kv$ has an
        effective dimension of one across the whole ensemble. Adding structure
        later means re-running.

  \item \textbf{Convective adjustment as a confound.} A large implicit
        diffusivity is applied wherever the water column is unstable. At the
        high-mixing end, explicit mixing may suppress convection, so part of the
        response would reflect convection triggering less often rather than a
        smooth dependence on $\kv$. Should convective activity be diagnosed per
        member?

  \item \textbf{Is ten years enough re-equilibration?} If the overturning has not
        settled, the sensitivities partly reflect drift rather than a new
        equilibrium.

  \item \textbf{Surrogate framing.} A global field-to-field map, or a local
        operator predicting sensitivity from local state variables? This decides
        whether forward diagnostics must be written at the 5-day adjoint cadence
        rather than monthly. Cheap to decide now, expensive to correct later.

  \item \textbf{Scope of the cost function.} All sensitivities here are
        $\partial J/\partial x$ for a single $J$, so a surrogate trained on them
        predicts that metric only. A general emulator needs $J$ to vary, and each
        additional cost function requires its own adjoint runs across the whole
        ensemble.
\end{enumerate}

Questions 4 and 5 are the ones most worth settling before the full ensemble.
Question 2 determines whether the ensemble can serve as surrogate training data
at all. Question 6 determines everything downstream.

%-----------------------------------------------------------------------
\section{Sequence and decision points}
\label{sec:sequence}
%-----------------------------------------------------------------------

\begin{center}
\begin{tabularx}{\textwidth}{@{}Y{0.2}Y{2.3}Y{0.5}@{}}
\toprule
\textbf{Step} & \textbf{Content} & \textbf{Cost}\\
\midrule
1 & Build and verify against the reference result; benchmark to confirm the
    adjoint/forward ratio and the gradient-check saving & $<$ 1 node-h\\
2 & Generate the seven diffusivity fields and configurations & ---\\
3 & Pilot member (8 $\times$ reference): 10-year leg from 2170; check overturning
    drift and convection at 2180 & $\sim$1.6 h\\
4 & Pilot adjoint 2180 $\rightarrow$ 2185; compare against the reference
  & $\sim$9 h\\
5 & \textbf{Review --- confirm the design before committing the remainder}
  & ---\\
6 & Remaining six members, run concurrently & $\sim$64 node-h\\
7 & Cross-member comparison; finite-difference validation of the adjoint gradient
  & ---\\
\bottomrule
\end{tabularx}
\end{center}

Steps 3 and 4 constitute one of seven runs, about 15\,\% of the ensemble, and
would expose a design problem before the remainder is committed. The strong
perturbation is chosen for the pilot because it gives the clearest signal on both
re-equilibration time and convective suppression.

Costs assume the gradient check is disabled for members, per
Section~\ref{sec:grdchk-cost}. At the reference run's settings the same steps
would cost about 1.6 h, 23 h and 150 node-h.

%=======================================================================

